\documentclass[reqno]{amsart}
\usepackage{a4wide}
\usepackage{hyperref}
\usepackage{amssymb}
\usepackage{graphicx} 

\newcommand{\R}{\ensuremath{\mathbb{R}}}
\newcommand{\Z}{\ensuremath{\mathbb{Z}}}
\newcommand{\G}{\ensuremath{\mathbb{H}}}
\newcommand{\N}{\ensuremath{\mathbb{N}}}
\newcommand{\K}{\ensuremath{\mathbb{K}}}
\newcommand{\PT}{\ensuremath{\mathbb{P}}}
\newcommand{\T}{\ensuremath{\mathbb{T}}}
\newcommand{\C}{\ensuremath{\mathbb{C}}}
\newcommand{\I}{\ensuremath{\mathbb{I}}}
\newcommand{\A}{\ensuremath{\mathbb{A}}}
\newcommand{\X}{\ensuremath{\mathbb{X}}}
\newcommand{\D}{\ensuremath{\mathbb{D}}}
\newcommand{\floor}[1]{\lfloor #1 \rfloor}

\newcommand{\Cerg}{\ensuremath{\mathbb{C}_{\textup{erg}}}}
\newcommand{\Chilg}{\ensuremath{\mathbb{C}_{\mathcal H}}}
\newcommand{\Rerg}{\ensuremath{\gamma}}
\newcommand{\Ierg}{\ensuremath{\omega}}

\DeclareMathOperator{\myRe}{Re}
\DeclareMathOperator{\myIm}{Im}


\AtBeginDocument{{\noindent\small
\emph{Electronic Journal of Differential Equations},
Vol. 2026 (2026), No. 04, pp. 1--20.\newline
ISSN: 1072-6691. URL: https://ejde.math.txstate.edu, DOI: 10.58997/ejde.2026.04}
\thanks{\copyright 2026. This work is licensed under a CC BY 4.0 license.}
\vspace{8mm}}

\begin{document}
\title[\hfilneg EJDE-2026/04\hfil Ergodic complex plane and cylinder transformation]
{Ergodic complex plane and cylinder transformation on a periodic time scale}

\author[J. M. Davis, B. J. Jackson, D. Poulsen \hfil EJDE-2026/04\hfilneg]
{John M. Davis, Billy J. Jackson, Dylan Poulsen}

\address{John M. Davis \newline
Department of Mathematics,
Baylor University, Waco, TX 76798, USA}
\email{John\_M\_Davis@baylor.edu}

\address{Billy J. Jackson \newline
Department of Mathematics,
Statistics, and Computer Science,
University of Illinois, Chicago, IL 60607, USA}
\email{bjjack06@uic.edu}

\address{Dylan Poulsen \newline
Department of Mathematics and Computer Science,
Washington College,
Chestertown, MD 21620, USA}
\email{dpoulsen2@washcoll.edu}

\thanks{Submitted November 10, 2025. Published January 12, 2026.}
\subjclass[2020]{93D05, 93D30, 37B25, 34N05, 26E70}
\keywords{Time scales; frequency; Hilger complex plane; stability}

\begin{abstract}
We introduce the ergodic complex plane, a global analogue of Hilger's local complex
plane on time scales, which simultaneously encodes exponential growth and frequency.
Averaging the cylinder transformation on a periodic time scale leads to the notions
of ergodic growth rate and ergodic frequency, unifying local and global stability
perspectives. This yields the ergodic cylinder transformation, a univalent map inducing
an orthogonal curvilinear coordinate system on the regressive complex plane. Within
this framework, we develop a decomposition analogous to Hilger's real and imaginary
parts, and define the box plus operation, extending the circle plus operation globally.
\end{abstract}

\maketitle
\numberwithin{equation}{section}
\newtheorem{theorem}{Theorem}[section]
\newtheorem{lemma}[theorem]{Lemma}
\newtheorem{definition}[theorem]{Definition}
\newtheorem{proposition}[theorem]{Proposition}
\newtheorem{remark}[theorem]{Remark}
\newtheorem{example}[theorem]{Example}
\allowdisplaybreaks

\section{Hilger's complex plane and exponential stability}

We begin with a brief summary of the Hilger complex plane \cite{BoPeD} on a time
scale $\mathbb{T}$. To match the narrative of later sections, we frame the Hilger
complex plane entirely in terms of the cylinder transformation, which is a different
perspective from \cite{BoPeD}.

Let $\mathbb{T}$ be a time scale and let $t \in \mathbb{T}$. Define the \emph{Hilger complex numbers}
at $t\in\mathbb{T}$, denoted $\mathbb{C}_{\mu(t)}$, by $\mathbb{C}_{\mu(t)} = \mathbb{C} \setminus \{-1/\mu(t)\}$. Defining $\Omega(\mu(t)) := \pi/\mu(t)$, let the \emph{cylinder strip} at $t\in\mathbb{T}$ be defined as
\[
\mathbb{C}^{\Omega(\mu(t))} := \{z \in \mathbb{C} \mid -\Omega(\mu(t)) < \myIm(z) \leq \Omega(\mu(t)) \}.
\]
\begin{definition} \rm
Let $\mathbb{T}$ be a time scale and let $t \in \mathbb{T}$. Define the
\emph{cylinder transformation} $\xi_{\mu(t)}:\mathbb{C}_{\mu(t)} \to \mathbb{C}^{\Omega(\mu(t))}$ by
 \begin{equation}\label{cyl}
 \xi_{\mu(t)}(z):=
 \begin{cases}
 \frac{\operatorname{Log}(1+z \mu(t))}{\mu(t)}, & \mu(t)>0,\\
 z, & \mu(t)=0.
 \end{cases}
 \end{equation}
\end{definition}

Note that $\xi_{\mu(t)}$ is one-to-one, and so $\xi_{\mu(t)}^{-1}$ exists.
In fact, it has the explicit form
\begin{equation}\label{invcyl}
 \xi_{\mu(t)}^{-1}(z):=
 \begin{cases}
 \frac{e^{\mu(t) z} -1}{\mu(t)}, & \mu(t)>0,\\
 z, & \mu(t)=0.
 \end{cases}
 \end{equation}

The \emph{Hilger real part} at $t\in\mathbb{T}$ of $z \in \mathbb{C}_{\mu(t)}$ can be defined as
 $$
 \myRe_{\mu(t)}(z)= \xi_{\mu(t)}^{-1}(\myRe(\xi_{\mu(t)}(z))).
 $$
The \emph{Hilger imaginary part} at $t \in \mathbb{T}$ of $z$ is given by
 $$
 \myIm_{\mu(t)}(z) = \myIm(\xi_{\mu(t)}(z)),
 $$
while the \emph{Hilger pure imaginary part} of $z$ at $t \in \mathbb{T}$ is given by
 $$
 \stackrel{o}{\iota} \myIm_{\mu(t)}(z) := \xi^{-1}_{\mu(t)}( i \myIm(\xi_{\mu(t)}(z))).
 $$
It is straightforward to establish the following identities (which are usually taken as definitions).

\begin{proposition} \label{lem:HilXi}
Let $\mathbb{T}$ be a time scale and let $t \in \mathbb{T}$. Then
\begin{gather*}
\myRe_{\mu(t)}(z)  = \lim_{\tau \searrow \mu(t)} \frac{|1+\tau z|-1}{\tau}, \\
\myIm_{\mu(t)}(z)  = \lim_{\tau \searrow \mu(t)}\frac{\operatorname{Arg}(z\tau+1)}{\tau}, \\
\stackrel{o}{\iota} \myIm_{\mu(t)}(z)
 = \lim_{\tau \searrow \mu(t)}\frac{e^{i \myIm_{\mu(t)}(z) \tau}-1}{\tau},
\end{gather*}
where $\operatorname{Arg}$ is the principal argument satisfying $-\pi < \operatorname{Arg}(z) \leq \pi$.
\end{proposition}

The \emph{circle plus operation} $\oplus_{\mu(t)}$ is defined by
$a \oplus_{\mu(t)} b:=a+b+\mu(t) ab$. The \emph{circle minus operation}
$\ominus_{\mu(t)}$ is defined to be the additive inverse of $\oplus_{\mu(t)}$.
Bohner and Peterson
\cite[Theorem~2.24]{BoPeD}  establish that $(\mathbb{C}_{\mu(t)}, \oplus_{\mu(t)})$
is homomorphic to $(\mathbb{C}^{\Omega(\mu(t))},+ (\text{mod } 2 \pi i/\mu(t)))$ with
group homomorphism $\xi_{\mu(t)}$. Although not explicitly stated in \cite{BoPeD},
it follows from the fact that
$\xi_{\mu(t)}$ is a bijection that
$(\mathbb{C}_{\mu(t)}, \oplus_{\mu(t)})$ is isomorphic to
$(\mathbb{C}^{\Omega(\mu(t))},+ (\text{mod } 2 \pi i/\mu) )$.

The group isomorphism $\xi_{\mu(t)}$ implies we can understand the circle plus
addition via, for $w,z \in \mathbb{C}_{\mu(t)}$,
\begin{equation}
\label{eq:oplusViaXi}
z \oplus_{\mu(t)} w = \xi_{\mu(t)}^{-1}(\xi_{\mu(t)}(z)
+ \xi_{\mu(t)}(w) \pmod{2 \pi i/\mu(t)}).
\end{equation}

Using circle plus, for fixed $t \in \mathbb{T}$, every $z \neq -\frac{1}{\mu(t)}$ has the
\emph{Hilger decomposition at time $t$} given by
$$
z=\text{Re}_{\mu(t)}(z) \oplus_{\mu(t)} \stackrel{o}{\iota}
\operatorname{Im}_{\mu(t)}(z).
$$

The Hilger complex plane, $\mathbb{C}_{\mu(t)}$, is shown in
Figure~\ref{fig:complexplane} for $\mu(t)\neq 0$.
Important components of $\mathbb{C}_{\mu(t)}$ are the \emph{Hilger disk}
at $t\in \mathbb{T}$, given by
$$
 \mathcal{H}_{\mu(t)} := \{ z \in \mathbb{C}_{\mu(t)} \mid \myRe_{\mu(t)}(z) < 0 \},
$$
which can be thought of as the preimage of the left half-plane of
$\mathbb{C}^{\Omega(\mu(t))}$ under $\xi_{\mu(t)}$.
The \emph{Hilger circle} at $t \in \mathbb{T}$ is the boundary of $\mathcal{H}_{\mu(t)}$,
which we denote by $\mathbb{I}_{\mu(t)}$.
The Hilger circle can be thought of as the preimage of the imaginary axis of
$\mathbb{C}^{\Omega(\mu(t))}$ under $\xi_{\mu(t)}$.

\begin{figure}[t]
  \centering
  \includegraphics[width=0.7\textwidth]{fig_complexplane.pdf}
  \caption{Hilger's complex plane, $\mathbb{C}_{\mathcal{H}}$. The $z$ inside the circle have negative
  Hilger real part, the $z$ on the circle have zero Hilger real part, and the
  $z$ outside the circle have positive Hilger real part. Points $z$ on the
  Hilger real axis $\mathbb{R}_{\mu(t)}$ (the solid ray on the real axis) are such
  that $e_{z}(t,t_0)>0$ for all $t$, while points on the Hilger alternating axis
  $\mathbb{A}_{\mu(t)}$ (the dotted ray on the real axis) are such that $e_{z}(t,t_0)$
  changes sign at each $t \in \mathbb{T}$.}
  \label{fig:complexplane}
\end{figure}



The cylinder transformation maps contours of equal Hilger imaginary part
(which are rays emanating from $-1/\mu(t)$) to horizontal lines in
$\mathbb{C}^{\Omega(\mu(t))}$. Also, contours of equal Hilger real part
(which are circles centered at $-1/\mu(t)$) are mapped to vertical lines in
$\mathbb{C}^{\Omega(\mu(t))}$. In particular, these contours form an orthogonal curvilinear coordinate system on $\mathbb{C}_{\mu(t)}$.

\begin{remark} \rm
The cylinder transformation is strongly related to the concepts of growth rate and
frequency on time scales with constant graininess. On the time scale
$\mathbb{T}=h\mathbb{Z}$, the contours of equal Hilger imaginary part correspond
to frequency given by
\begin{equation}
\label{eq:usefulIm}
\omega_h(z):=\operatorname{Im}_h(z) = \myIm(\xi_h(z)) = \frac{\operatorname{Arg}(1+z h)}{h}.
\end{equation}
This is the same quantity that represents frequency in Hilger's definition of the
time scale sinusoids \cite{HiA}. Moreover, the contours of equal Hilger real part
all have the same exponential growth rate corresponding to the Lyapunov exponent
\begin{equation}
\label{eq:usefulRe}
\gamma_{h}(z):= \frac{\ln(1+h\, \text{Re}_h(z))}{h}
= \xi_{h}(\myRe_h(z)) = \myRe(\xi_h(z)) = \frac{\ln|1+hz|}{h}.
\end{equation}
\end{remark}
Similarly, on the time scale $\mathbb{T}=\mathbb{R}$, $\xi_0(z) = z$, and the contours
of equal exponential growth rate are vertical lines determined by the Lyapunov
exponent $\gamma_0(z):=\text{Re}(z)$, while the contours of equal frequency are
horizontal lines given by $\omega_0(z)=\operatorname{Im}(z)$.

We say that $p:\mathbb{T} \to \mathbb{R}$ is \emph{regressive} if $1+\mu(t)p(t)\not=0$
for all $t\in\mathbb{T}$.
For regressive $p$, the \emph{time scale exponential function} $e_p(t,s)$ is defined by
 \[
 e_{p}(t,s) := \exp\Big(\int_{s}^{t} \xi_{\mu(t)}(p(\tau)) \Delta \tau \Big).
 \]
Following Karpuz \cite{Ka}, denote the set of all \emph{regressive points} in the
complex plane by
$$
\mathbb{C}_{\mu(\mathbb{T})}:=\{z \in \mathbb{C} \mid 1+\mu(t) z \neq 0 \text{ for all } t \in \mathbb{T}\}.
$$

The definitions related to Hilger's complex plane are \emph{local} and \emph{dynamic}
in nature; that is, they depend on and change in $t$. In 2003, P\"{o}tzsche, Siegmund,
 and Wirth introduced the concept of a \emph{global} and \emph{static} exponential
 growth rate for the time scale exponential function.

\begin{theorem}[P\"{o}tzsche, Siegmund, and Wirth \cite{PoSiWi}]\label{PSW1}
Let $\mathbb{T}$ be a time scale which is unbounded above and let $\lambda
\in \mathbb{C}$. Then the scalar equation
 $$
 x^{\Delta}(t)=\lambda x(t), \quad x(t_0)=x_0,
 $$
is exponentially stable if and
only if one of the following conditions is satisfied for arbitrary
$t_0 \in \mathbb{T}$:
\begin{itemize}
\item[(C1)] $\gamma(\lambda):=  \limsup_{T \to
\infty}\frac{1}{T-t_0}\int_{t_0}^{T}\lim_{s \searrow
\mu(t)} \frac{\ln|1+\lambda s|}{s}\Delta t <0$,

\item[(C2)] For every
$T \in \mathbb{T}$, there exists a $t \in \mathbb{T}$ with $t>T$ such that
$1+\mu(t)\lambda=0$,
\end{itemize}
where we use the convention $\ln 0=-\infty$ in \textup{(C1)}.
\end{theorem}

Note that condition (C1) can be restated using \eqref{eq:usefulRe} as
\begin{itemize}
\item[(C1*)] $ \limsup_{T \to
\infty}\frac{1}{T-t_0}\int_{t_0}^{T}\lim_{s \searrow
\mu(t)} \gamma_s(\lambda) \Delta t <0$.
\end{itemize}
This theorem naturally leads to the following definition.

\begin{definition}[\cite{PoSiWi}] \rm
Given a time scale $\mathbb{T}$ which is unbounded above, for
arbitrary $t_0 \in \mathbb{T}$, define the sets
 \begin{gather*}
 \mathcal{S}_{\mathbb{C}}(\mathbb{T}):=\{\lambda \in \mathbb{C}: \text{ (C1) holds}\},\\
 \mathcal{S}_{\mathbb{R}}(\mathbb{T}):=\{\lambda \in \mathbb{R}: \text{ (C2) holds}\}.
 \end{gather*}
Then the \emph{set of exponential stability} for $\mathbb{T}$ is given by
$$
\mathcal{S}(\mathbb{T}):=\mathcal{S}_{\mathbb{C}}(\mathbb{T}) \cup \mathcal{S}_{\mathbb{R}}(\mathbb{T}).
$$
\end{definition}

The calculation of the set of exponential stability is simplified for periodic
time scales [\cite[Lemma~3.4(b)]{PoSiWi}, , which will be a focus of this paper.

\begin{definition} \rm
 A time scale $\mathbb{T}$ is \emph{periodic} if there exists $L>0$ such that
 $t \in \mathbb{T}$ implies $t+L \in \mathbb{T}$ and $\mu(t+L)=\mu(t)$ for all $t \in \mathbb{T}$.
We call $L$ a \emph{period} of the time scale.
\end{definition}


\begin{remark} \rm
Throughout this paper, we will illustrate various results on the following
prototypical time scales.
 \begin{enumerate}
 \item For $\mu_1,\mu_2,\dots,\mu_n>0$, set $L=\mu_1+\mu_2+\cdots+\mu_n$.
 Then $\mathbb{T}_{\mu_1,\mu_2,\dots,\mu_n}$ is the time scale starting at 0 whose
 graininesses form the periodic sequence $\{\mu_1,\mu_2,\dots,\mu_n\}$, i.e.,
 $$
 \mathbb{T}_{\mu_1,\mu_2,\dots,\mu_n}
 =\Big\{0,\mu_1,\mu_1+\mu_2,\dots,\underbrace{\mu_1+\cdots+\mu_n}_{L},
 L+\mu_1,L+\mu_1+\mu_2,\dots,2L,\dots\Big\}.
 $$

 \item $\mathbb{P}_{a,b}$ is the periodic time scale starting at 0 with an
 interval of length $a$ followed by a gap of length $b$, i.e.,
 $$
 \mathbb{P}_{a,b}=[0,a]\cup[a+b,2a+b]\cup[2a+2b,3a+2b]\cup\cdots
 $$
\end{enumerate}
\end{remark}

A key insight for this paper is that since the cylinder transformation is linked to both
growth rate (via the Hilger real part) and frequency (via the Hilger imaginary part) on
constant graininess time scales, and since the average of the real part of the cylinder
transformation is related to global exponential growth rate, it seems reasonable that
the average of the imaginary part of the cylinder transformation is related to the
global frequency. 

\section{Ergodic growth rate and frequency}

To solidify the connection between the average of the cylinder transformation and the
growth rate and frequency of the time scale exponential, notice for $\lambda \in \mathbb{C}_h$
that the time scale exponential function decomposes as
\begin{equation} \label{eq:exp_gamma_omega}
\begin{aligned}
 e_{\lambda}(t,t_0)
 &=  \exp\Big(\int_{t_0}^{t} \xi_{\mu(\tau)}(\lambda) \Delta \tau \Big) \\
 &= \exp\Big( \frac{\int_{t_0}^{t} \xi_{\mu(\tau)}(\lambda) \Delta \tau}
  {t-t_0} (t-t_0) \Big) \\
 &= \exp\Big( \frac{\int_{t_0}^{t} \myRe(\xi_{\mu(\tau)}(\lambda))
  \Delta \tau}{t-t_0} (t-t_0) \Big)
  \exp\Big( i \frac{\int_{t_0}^{t} \myIm(\xi_{\mu(\tau)}(\lambda)) \Delta \tau}{t-t_0}
   (t-t_0) \Big) \\
 &\overset{\eqref{eq:usefulIm},\eqref{eq:usefulRe}}{=}
 \exp\Big( \frac{\int_{t_0}^{t} \gamma_{\mu(\tau)}(\lambda) \Delta \tau}{t-t_0}
  (t-t_0) \Big)
  \exp\Big( i \frac{\int_{t_0}^{t} \omega_{\mu(\tau)}(\lambda) \Delta \tau}{t-t_0}
  (t-t_0) \Big) \\
 &=  \exp\Big( \frac{\int_{t_0}^{t} \gamma_{\mu(\tau)}(\lambda) \Delta \tau}{t-t_0}
  (t-t_0) \Big)   \\
 &\quad \times \Big[\cos\Big( \frac{\int_{t_0}^{t} \omega_{\mu(\tau)}(\lambda)
 \Delta \tau}{t-t_0} (t-t_0) \Big) + i \sin\Big( \frac{\int_{t_0}^{t}
 \omega_{\mu(\tau)}(\lambda) \Delta \tau}{t-t_0} (t-t_0) \Big) \Big].
 \end{aligned}
\end{equation}
Therefore, over an interval $[t,t_0]_{\mathbb{T}}$, the growth rate and frequency of the
time scale exponential function are averages (in the integral sense) of
$\gamma_{\mu(\tau)}$ and $\omega_{\mu(\tau)}$, respectively.

Since we are averaging the cylinder transformation, which is a logarithm for
$\mu(t)>0$, we want to be careful about the branch cut. We begin by defining a subset
of the complex plane which avoids the branch cuts of \eqref{cyl} for all $\tau\in\mathbb{T}$.
Assuming \eqref{P}, $\mu_{\max}:=\max \{ \mu(t) \mid t \in \mathbb{T} \}$ exists, so define
\[
 B:= (-\infty, -1/\mu_{\max}],\quad
 B^c= \mathbb{C} \setminus B.
\]
Motivated by the form \eqref{eq:exp_gamma_omega}, we now formalize the averaged
quantities introduced above.

\begin{definition}
\label{def_gamma_omega} \rm
Let $\mathbb{T}$ be a time scale and $x,y \in \mathbb{R}$ such that $x+iy \in B^c$.
The \emph{ergodic growth rate of $e_{x+iy}(t,t_0)$ on $[t_0,t]_\mathbb{T}$ off the branch cut}
is given by
\begin{equation}
 \label{eq:gammatt0}
 \gamma_{B^c}(x,y,t_0,t) :=\frac{1}{t-t_0}\int_{t_0}^{t} \gamma_{\mu(\tau)}(x+iy)
 \Delta \tau
\end{equation}
Similarly, the \emph{ergodic frequency of $e_{x+iy}(t,t_0)$ on $[t_0,t]_\mathbb{T}$ off the
branch cut} is given by
 \begin{equation}
 \omega_{B^c}(x,y,t_0,t) := \frac{1}{t-t_0}\int_{t_0}^{t}
 \omega_{\mu(\tau)}(x+iy)\,\Delta \tau.\label{w}
 \end{equation}
\end{definition}

\begin{remark} \rm
Several remarks are in order regarding $\gamma_{B^c}(x,y,t_0,t)$ and
$\omega_{B^c}(x,y,t_0,t)$.

(1) We use the term \emph{ergodic} to describe the growth rate and frequency here
because, if one considers the space $\mathbb{X}=\mu^*(\mathbb{T})$ of all extended graininesses
of $\mathbb{T}$ \cite{JaDaE} in the order that they appear in $\mathbb{T}$, then the global (spatial)
averages of the growth rates and frequencies derived from $\mathbb{X}$ are the
pointwise limit of the local (time) averages of the local growth rates and frequencies
derived from $\mathbb{X}$ as defined above. We may treat the collection of asymptotic
extended graininesses as an irreducible Markov chain, which guarantees the ergodicity
of the associated stochastic process. See \cite{PoDaGr}.

(2) By \eqref{eq:usefulRe}, $\gamma_{h}(x+iy)$ has a few useful forms. The choice of
form gives different interpretations to the equation. Writing
$\gamma_{h}(x+iy) = \xi_{h}(\myRe_{h}(x+iy))$ is useful when a characterization of
growth rates in terms of the Hilger real part of $x+iy$ is advantageous.
By contrast, writing $\gamma_{h}(x+iy) = \myRe(\xi_{h}(x+iy))$ is in line with
P\"otzsche, Siegmund, and Wirth's condition.

(3) Using properties of the logarithm, we get
$$
\gamma_{h}(x+iy) = \ln(1+h \myRe_{h}(x+iy))^{1/h}.
$$
This implies $\gamma_{B^c}(x,y,t_0,t)$ is the natural log of the geometric mean (in
the time scales sense) on $[t_0,t]_\mathbb{T}$ of
\begin{equation}\label{loc}
 \lim_{s \searrow\mu(t)}(1+\operatorname{Re}_{s}(\lambda) s)^{1/s}
 = \lim_{s \searrow\mu(t)}|1+\lambda s|^{1/s}.
\end{equation}
Thus, \eqref{loc} can be interpreted as a type of local growth rate of $e_{x+iy}(t,t_0)$
at time $t$. Since the geometric mean is the best measure of average local growth rates,
expressing $\gamma_{h}(x+iy)$ in this way shows that \eqref{eq:gammatt0} is an effective
 measure of the average growth rate of the time scale exponential function over
$[t_0,t]_{\mathbb{T}}$. This observation was made earlier in \cite{DaGrMaJa,JaDaE} --albeit
from the geometric perspective-- by geometrically averaging Hilger circles.

(4) The formulation in \eqref{eq:exp_gamma_omega} is similar to the one provided by
 Karpuz \cite{Ka}:
 \[% \label{karpuz}
 e_{z}(t,t_0) = e_{\text{Re}_{\mu(t)}(z)}(t,t_0)
 \Big[ \cos\Big(\int_{t_0}^{t} \operatorname{Im}_{\mu(\eta)}(z) \Delta \eta\Big)
 + i \sin\Big(\int_{t_0}^{t} \operatorname{Im}_{\mu(\eta)}(z) \Delta \eta\Big) \Big].
 \]
To use the terminology of \cite{JaDalap}, Karpuz gave a type II (time scale)
exponential representation whereas \eqref{eq:exp_gamma_omega} is a type I (continuous)
exponential representation.

(5) In this new notation, condition (C1*) becomes
 \begin{equation} \label{supgamma}
 \limsup_{T \to \infty} \gamma_{B^c}(x,y,t_0,T) < 0,
 \end{equation}
and thus (C1*) is indeed a requirement that the largest cluster point of the average
exponential growth rate over the tail of the time scale is negative.
Therefore, \eqref{supgamma} is a condition on the \emph{global} exponential growth rate.
 Moreover, the formulation here ties $\gamma_{B^c}(x,y,t_0,t)$ to the local Lyapunov
 exponent $\gamma_{\mu(t)}$, in the spirit of P\"otzsche, Siegmund, and Wirth's
 (C1) condition.

(6) Note that the dynamic equation
\begin{equation}
z^{\Delta}(t)=\lambda z(t), \quad z(t_0)=z_0,
\label{expequation}
\end{equation}
for $\lambda \in B^c$ has a growth rate of $\gamma_{B^c}(x,y,t_0,\infty)$ rather than a growth rate of $\lambda$. Furthermore, unlike on $\mathbb{R}$ where the terms \emph{growth rate} and \emph{exponential order} are often used synonymously, these two quantities will be distinct in general. Indeed, if $\lambda$ in \eqref{expequation} is real and positive, then $\lambda$ is best described as the \emph{exponential order} of the equation.
 \end{remark}

We now see that we can understand the \emph{global} growth rate and frequency of the
time scale exponential function by understanding the asymptotic behavior of the
ergodic growth rate and ergodic frequency as $t\to\infty$.
 Since this asymptotic behavior is complicated in general (indeed, the limit as
$t\to\infty$ may not exist), we begin by studying the asymptotic behavior with the
following simplifying assumption:
\begin{equation}
\mathbb{T} \text{ is a periodic time scale with a period of } L>0
\text{ and }0 \in \mathbb{T}.
\tag{P}
\label{P}
\end{equation}
Under this assumption, we have
\[
\lim_{t \to \infty} \gamma_{B^c}(x,y,t,0) = \gamma_{B^c}(x,y,L,0),
\]
and similarly
\[
\lim_{t \to \infty} \omega_{B^c}(x,y,t,0) = \omega_{B^c}(x,y,L,0).
\]
Therefore, we can simplify notation by rewriting \eqref{eq:gammatt0}, \eqref{w},
for $x+iy \in B^c$, as
\begin{gather}
 \gamma_{B^c}(x,y) :=\frac{1}{L}\int_0^{L} \lim_{s \searrow
\mu(\tau)} \gamma_s(x+iy)\,\Delta \tau, \label{g21} \\
\omega_{B^c}(x,y) := \frac{1}{L}\int_{0}^{L} \lim_{s \searrow
\mu(\tau)} \omega_s(x+iy)\,\Delta \tau.\label{w1}
\end{gather}

\begin{remark} \rm
While the assumption of periodicity may appear to be limiting, a larger class of
aperiodic time scales can be analyzed using periodic techniques.
We will call a time scale \emph{simple} if there is a representative periodic time
scale which has the same limiting functions $\gamma_{B^c}$ and $\omega_{B^c}$.

To illustrate this, consider the symmetric time scale
$\mathbb{T}_{\text{PTM}} =\{0,t_1,t_2,\ldots\}$ with the graininess function $\mu(t_n)$
 equal to the $n^{\text{th}}$ term of the celebrated Prouhet-Thue-Morse sequence on
the symbols one and two \cite{allouche1999ubiquitous}.
Since the Prouhet-Thue-Morse sequence is aperiodic, $\mathbb{T}_{\text{PTM}}$ is also aperiodic.
Additionally, since the symbols of the sequence appear with equal weight in the limit
as the sequence length increases,
\begin{align*}
 \limsup_{T \to \infty} \gamma_{B^c}(x,y,t_0,T)
&= \limsup_{T \to \infty} \frac{1}{T-t_0} \int_{t_0}^{T} \frac{\ln|1+(x+iy) s|}{s} \Delta \tau\\
&= \frac{1}{3} \left[ \ln|1+(x+iy)| + \ln|1+2(x+iy)| \right].
\end{align*}
This last expression is equal to $\gamma_{B^c}(x,y)$ on the periodic time scale
$\mathbb{T}_{1,2}$. Similar results hold for calculating the limiting value of $\omega_{B^c}$.

In this paper, any result which assumes \eqref{P} can be extended to a simple time
scale via its representative (asymptotically equivalent) periodic time scale since our
arguments rely only on the limiting values of $\gamma_{B^c}$ and $\omega_{B^c}$
rather than the precise nature of $\mathbb{T}$ itself.
\end{remark}

In light of the definitions above, the level curves of $\gamma_{B^c}(x,y)$ are the
curves along which $e_{x+iy}(t,t_0)$ has a constant global exponential growth rate.
The level curves of $\omega_{B^c}(x,y)$ are the curves along which $e_{x+iy}(t,t_0)$
has a constant global (signal) frequency. As we will see, these two families of level
curves induce a natural coordinate system on $\mathbb{C}_{\mu(\mathbb{T})}$.

Using these definitions, we can reformulate $\mathcal{S}_{\mathbb{C}}(\mathbb{T})$ as
\begin{equation}\label{pswtype}
 \mathcal{S}_{\mathbb{C}}(\mathbb{T}) = \{x+iy \in \mathbb{C} \mid \gamma_{B^c}(x,y) <0\}.
\end{equation}
Equation \eqref{pswtype} is reminiscent of condition (C1) and $\mathcal{S}_\mathbb{C}(\mathbb{T})$
has been thought of as an average of Hilger circles \cite{DaGrMaJa}.

Finally, since $\gamma_{B^c}$ and $\omega_{B^c}$ are the average value of
$\myRe(\xi_{\mu(\tau)})$ and $\myIm(\xi_{\mu(\tau)})$, respectively, it makes sense to
study directly the complex function given by the average of $\xi_{\mu(\tau)}$.
In the next section, we will define and study the properties of a map $\overline{\xi}$,
called the \emph{ergodic cylinder transformation}, from the $xy$ plane to the
$\gamma \omega$ plane, i.e., a mapping from the pair $(x,y)$ to the pair
$(\gamma(x,y),\omega(x,y))$ for periodic time scales.

\section{The ergodic cylinder transformation}

We begin by defining the codomain of the ergodic cylinder transformation. We define
 \begin{gather*}
 \Omega := \sup_{x+iy \in B^c} \omega_{B^c}(x,y),\\
 \mathbb{C}^{\Omega} := \{z \in \mathbb{C} \mid -\Omega< \myIm(z) \leq \Omega\}.
 \end{gather*}
If $\Omega=\infty$, then $\mathbb{C}^{\Omega} = \mathbb{C}$.

\begin{definition}
\label{def:xibarBc} \rm
Assume \eqref{P}. Suppose $x+iy\in B^c$. We define the \emph{ergodic cylinder
 transformation off the branch cut}, $\overline{\xi}_{B^c}: B^c \to \mathbb{C}^{\Omega}$, by
$$
 \overline{\xi}_{B^c}(x+iy):=\frac{1}{L}\int_{0}^{L} \xi_{\mu(t)}(x+iy) \Delta t.
$$
\end{definition}

It is straightforward to show from \eqref{eq:exp_gamma_omega}, for $x+iy \in B^c$,
\begin{equation} \label{eq:gamOm}
 \overline{\xi}_{B^c}(x+iy) = \gamma_{B^c}(x,y) + i\,\omega_{B^c}(x,y).
\end{equation}

Our goal is to extend $\overline{\xi}_{B^c}$ to a map $\overline{\xi}:\mathbb{C}_{\mu(\mathbb{T})} \to \mathbb{C}^{\Omega}$
that induces an orthogonal curvilinear coordinate system on $\mathbb{C}_{\mu(\mathbb{T})}$ for which
the contours are given by the level curves of $\gamma$ and $\omega$. We will first show
that this is the case for $\overline{\xi}_{B^c}$, and then we will define the extension.

Analytic functions that are globally injective, or univalent, induce an orthogonal
coordinate system as the preimage of the rectangular coordinate system in the complex
plane via the Cauchy-Riemann equations. Therefore, we begin by showing
$\overline{\xi}_{B^c}:B^c \to \mathbb{C}^{\Omega}$ is analytic and univalent. To do so, we need the
following result, in the form first presented in \cite[Theorem~8]{AvAk}.

\begin{theorem}[Noshiro-Warschawski criterion]
If $f$ is analytic and nonconstant in a convex domain $D$, and
$$
\operatorname{Re}(e^{i \alpha} f'(z))>0
$$
for all $z \in D$ and for $\alpha \in \mathbb{R}$ fixed, then $f$ is univalent in $D$.
\end{theorem}

\begin{theorem}
\label{thm:analytic}
Suppose \eqref{P} holds. Then
\begin{enumerate}
 \item $\overline{\xi}_{B^c}:B^c \to \mathbb{C}^{\Omega}$ is analytic, and
 \item $\overline{\xi}_{B^c}:B^c \to \mathbb{C}^{\Omega}$ is univalent.
\end{enumerate}
\end{theorem}

\begin{proof}
Let $z=x+iy\in B^c$.

(1) Since the principal logarithm is analytic, $e_{z}$ is analytic for
$z \in \mathbb{C}_{\mu(\mathbb{T})}$ \cite{Ka}, and the nonregressive points lie in $B$,
it follows that $\overline{\xi}_{B^c}$ is analytic.

(2) To show univalence, from \cite{Ka}, the function $m_z:\mathbb{C}_{\mu(\mathbb{T})} \to \mathbb{C}$,
$$
m_z(s,t) = \int_{s}^{t} \frac{\Delta \tau}{1+z\mu(\tau)},
$$
is analytic in $z$, so
\begin{equation}
\begin{aligned}
\overline{\xi}_{B^c}'(z)
& = \frac{1}{L} \int_{0}^{L} \frac{\Delta \tau}{1+z \mu(\tau)}  \\
& = \frac{1}{L} \int_{0}^{L} \frac{(1+x \mu(\tau))\Delta \tau}{|1+(x+iy)\mu(\tau)|^2}
- \frac{i}{L} \int_{0}^{L} \frac{y \mu(\tau))\Delta \tau}{|1+(x+iy)\mu(\tau)|^2}.
\end{aligned} \label{calc}
\end{equation}

If $\mathbb{T}=\mathbb{R}$, univalence follows immediately since $\overline{\xi}_{B^c}(z) = z$.
If $\mathbb{T}\not=\mathbb{R}$, we consider three cases: $y<0,y=0$, and $y>0$. If $y>0$, let
$U := \{x+iy \in \mathbb{C} \mid y>0 \}$. Then, for $x+iy \in U$,
 \[
 \myRe(i\,\overline{\xi}_{B^c}'(x+iy))
  = \frac{1}{L} \int_{0}^{L} \frac{y \mu(\tau)\Delta \tau}{|1+(x+iy)\mu(\tau)|^2}>0,
 \]
since the integrand is nonnegative and not identically zero over $[0,L]_{\mathbb{T}}$ due to
the fact that $\mu(t) \not \equiv 0$, $y>0$, and (crucially) that
 $$
 \inf_{\tau \in [0,L]_{\mathbb{T}}} \frac{1}{|1+(x+iy) \mu(\tau)|^2} >0,
 $$
by \cite[Lemma 2.1]{Ka}. Since $U$ is convex, by the Noshiro-Warschawski criterion
(with $\alpha=\pi/2$), $\overline{\xi}_{B^c}$ is univalent on $U$.

Similarly, if $y<0$, $\overline{\xi}_{B^c}$ is univalent on the lower half-plane.
Note that if $z_1, z_2 \in \mathbb{C}$ with $\myIm(z_1)>0$ while $\myIm(z_2)<0$, then
$\overline{\xi}_{B^c}(z_1) \neq \overline{\xi}_{B^c}(z_2)$ since $\omega_{B^c}(x,y)>0$ for $y>0$
and $\omega_{B^c}(x,y)<0$ for $y<0$.

 Finally, if $y=0$ and $x> -1/\mu_{\max}$, $\omega_{B^c}(x,0) =0$ while
$\overline{\xi}_{B^c}'(x)>0$, which implies $\overline{\xi}_{B^c}$ is one-to-one for $x>-1/\mu_{\max}$.
Since $\omega_{B^c}$ is positive in the upper half-plane, negative in the lower-half
plane, and zero on the ray $(-1/\mu_{\max},\infty)$, and since $\overline{\xi}_{B^c}$ is
univalent on each of these spaces, it follows that $\overline{\xi}_{B^c}$ is globally
univalent.
\end{proof}

We now show that $\overline{\xi}_{B^c}$ can be extended to a map $\overline{\xi}:\mathbb{C}_{\mu(\mathbb{T})} \to \mathbb{C}$
which is globally univalent. We will call $\overline{\xi}$ the \emph{ergodic cylinder
transformation}. To do so, we treat the behavior along the branch cut $B$ carefully.
Since the branch cut consists of intervals between nonregressive points, and since
the properties of these intervals depend on properties of the time scale, we establish
some notation and technical lemmas to aid in the proof.

Let $M=\{\mu_n\}$ be the collection of distinct, positive graininesses of $\mathbb{T}\not=\mathbb{R}$,
ordered as
$$
\mu_{\max}=\mu_1>\mu_2>\cdots.
$$
Let $\mu_*=\inf_{n} \mu_n$ so that
$$
 -\frac{1}{\mu_*}<\cdots < -\frac{1}{\mu_2} < -\frac{1}{\mu_1}<0.
$$
Let
 \begin{gather*}
 I_n:=\Big(-\frac{1}{\mu_{n+1}},-\frac{1}{\mu_n}\Big),\ n=1,2,\dots,\\
 I_0:=
 \begin{cases}
 (-\infty,-1/\mu_*), & \mu_*>0,\\
 \emptyset, & \mu_*=0.
 \end{cases}
 \end{gather*}

\begin{lemma}
\label{lem:B_structure}
Assume \eqref{P}. Then $M$ is either finite or countably infinite. Moreover,
\begin{enumerate}
 \item If $M$ is finite,
 $$
 B\setminus \Big\{-\frac{1}{\mu_n}\Big\} = I_0 \cup \left(\bigcup_{n=1}^{N-1} I_n\right).
 $$

 \item If $M$ is countably infinite,
 $$
 B\setminus \Big\{-\frac{1}{\mu_n}\Big\} = \bigcup_n^\infty I_n.
 $$
\end{enumerate}
\end{lemma}

\begin{proof}
An arbitrary time scale has at most countably many graininesses.
 \begin{enumerate}
 \item If $M$ is finite, then $\mu_*>0$ so $I_0\not=\emptyset$ and the result follows.
 \item If $M$ is countably infinite, then we still have $ L=\sum_n \mu_n<\infty$,
 so it must be that $\mu_*=0$. Hence, $I_0=\emptyset$ and the result follows.
 \end{enumerate}
\end{proof}

\begin{theorem}
\label{thm:GlobalUnivalence}
$\overline{\xi}_{B^c}:B^c \to \mathbb{C}^{\Omega}$ can be extended to a globally univalent function
$\overline{\xi}:\mathbb{C}_{\mu(\mathbb{T})} \to \mathbb{C}^{\Omega}$.
\end{theorem}

\begin{proof}
We present the proof in two cases.
\smallskip

\noindent\textbf{Case 1: $M$ finite.}
Since $M$ is finite, we can list the $N$ distinct graininesses in order as
 \[
 \mu_{\max}=\mu_1>\mu_2>\mu_3>\cdots>\mu_N,
 \]
so that the corresponding non-regressive points satisfy
 \[
 -\frac{1}{\mu_N} < \cdots < -\frac{1}{\mu_2} < -\frac{1}{\mu_1} < 0.
 \]
Then, by Lemma \ref{lem:B_structure}, we have
$$
 B\setminus \Big\{-\frac{1}{\mu_n}\Big\}
 = I_0 \cup \left(\bigcup_{n=1}^{N-1} I_n\right).
$$
To extend $\overline{\xi}_{B^c}$, we consider
\begin{equation}\label{setup}
 \overline{\xi}_{B^c}^+(x):=\lim_{y\to 0^+}\overline{\xi}_{B^c}(x+iy),
 \quad\text{for } x\in B\setminus\{-1/\mu_n\}.
\end{equation}
Taking real and imaginary parts of $\overline{\xi}_{B^c}^+(x)$ gives
\begin{gather*}
 \gamma_B(x):=\operatorname{Re} \overline{\xi}_{B^c}^+(x)
 =\frac{1}{L}\int_{\{t\in[0,L]:\mu(t)=0\}} x\,\Delta t
  + \frac{1}{L}\int_{\{t\in[0,L]:\mu(t)>0\}} \frac{\ln|1+\mu(t) x|}{\mu(t)}
  \Delta t,
\\
\begin{aligned}
 \omega_B(x):=\operatorname{Im} \overline{\xi}_{B^c}^+(x)
 &=\frac{\pi}{L}\int_{\{t\in[0,L]:\mu(t)>0\}}
  \frac{\mathbf 1_{\{x<-1/\mu(t)\}}}{\mu(t)}\,\Delta t\\
 &=\frac{\pi}{L} \sum_{\{t\in[0,L]:\mu(t)>0\}} \mathbf 1_{\{x<-1/\mu(t)\}}\\
 &=\frac{\pi}{L}
 \begin{cases}
 \sum_{\{t\in[0,L]:\mu(t)\ge \mu_n\}} 1, & x\in I_n,\\
 \sum_{\{t\in[0,L]:\mu(t)>0\}} 1, & x\in I_0,
 \end{cases}
\end{aligned}
\end{gather*}
where $\mathbf{1}_A$ denotes the characteristic function on the set $A$.
Thus, $\omega_B$ is constant on each $I_n$; write $\omega_n:=\omega_B(x)$ for any
$x\in I_n$. Also, $\omega_B$ is constant on $I_0$ as well, so we let
$\omega_0:=\omega_B(x)$ for $x\in I_0$. Then
 $$
 \omega_1<\omega_2<\cdots<\omega_{N-1}<\omega_0.
 $$

For $x\in B\setminus\{-1/\mu_n\}$, differentiating under the integral yields
\begin{gather*}
\begin{aligned}
 \gamma_B'(x)
 &=\frac{1}{L}\int_{\{[0,L]:\mu(t)=0\}} 1\,\Delta t
 + \frac{1}{L}\int_{\{[0,L]:\mu(t)>0\}} \frac{1}{1+\mu(t) x}\,\Delta t\\
 &=p_0+\frac{1}{L}\int_{\{[0,L]:\mu(t)>0\}} \frac{1}{1+\mu(t) x}\,\Delta t,
\end{aligned}\\
\gamma_B''(x)
 =-\frac{1}{L}\int_{\{[0,L]:\mu(t)>0\}} \frac{\mu(t)}{(1+\mu(t) x)^2}\,\Delta t\ <\ 0,
\end{gather*}
where $$p_0:=\frac{1}{L}\int_{\{[0,L]:\mu(t)=0\}} 1\,\Delta t$$ is a constant.
Thus, $\gamma_B'$ is strictly decreasing on each $I_n$, $n=1,2,\dots$.
Moreover, on $I_n$, we have the end behavior
 \[
 \lim_{x\downarrow -1/\mu_{n+1}}\gamma_B'(x)=+\infty,\quad
 \lim_{x\uparrow -1/\mu_{n}}\gamma_B'(x)=-\infty.
 \]
Hence, there exists a unique $\alpha_n\in I_n$ with $\gamma_B'(\alpha_n)=0$.
Consequently,
 \[
 \gamma_B'(x)>0\ \text{on}\ \big(-\tfrac{1}{\mu_{n+1}},\,\alpha_n\big),
 \quad
 \gamma_B'(x)<0\ \text{on}\ \big(\alpha_n,\,-\tfrac{1}{\mu_n}\big).
 \]
We can similarly analyze the behavior of $\gamma_B$ on $I_0$. If $p_0=0$,
then $\gamma_B'(x)<0$ for $x\in I_0$. If $p_0>0$, then there exists a unique
$\alpha_0\in (-\infty,-1/\mu_N)$ such that $\gamma_B'(\alpha_0)=0$ and $\gamma_B'(x)>0$
for $x\in(-\infty,\alpha_0)$ while $\gamma_B'(x)<0$ for $x\in (\alpha_0,-1/\mu_N)$.

We define $\overline{\xi}_B:B\setminus\{-1/\mu_n\}\to\mathbb{C}$ by
 \begin{equation}\label{eff}
 \overline{\xi}_B(x)=
 \begin{cases}
 \gamma_B(x)+ i\,\omega_0, & p_0>0,\ x\in(-\infty,\alpha_0],\\[2pt]
 \gamma_B(x)- i\,\omega_0, & p_0>0,\ x\in(\alpha_0,-1/\mu_N),\\[2pt]
 \gamma_B(x)+ i\,\omega_0, & p_0=0,\ x\in(-\infty,-1/\mu_N),\\[2pt]
 \gamma_0(x)- i\,\omega_n, & x\in\big(-\tfrac{1}{\mu_{n+1}},\,\alpha_n\big],\;
  n=1,\dots,N-1,\\[2pt]
 \gamma_B(x)+ i\,\omega_n, & x\in\big(\alpha_n,\,-\tfrac{1}{\mu_n}\big),\ n=1,\dots,N-1,
 \end{cases}
 \end{equation}
Note that on each $I_n$, we use $-\overline{\xi}_{B^c}^+(x)$ for $x\in(-1/\mu_{n+1},\alpha_n]$,
but use $\overline{\xi}_{B^c}^+(x)$ for $x\in(\alpha_n,-1/\mu_n)$. As a preference, on $I_0$,
we use $\overline{\xi}_{B^c}^+(x)$ for $x\in(-\infty,\alpha_0]$, but use $-\overline{\xi}_{B^c}^+(x)$
for $x\in(\alpha_0,-1/\mu_N)$.

Finally, define $\overline{\xi}:\mathbb{C}_{\mu(\mathbb{T})} \to \mathbb{C}^{\Omega}$ by
 $$
 \overline{\xi}(z)=
 \begin{cases}
 \overline{\xi}_{B^c}(z), & z\in B^c,\\[2pt]
 \overline{\xi}_B(x), & z=x\in B\setminus\{-1/\mu_n\}.
 \end{cases}
 $$
We now show that $\overline{\xi}$ is globally univalent. We consider three subcases.
\smallskip

\noindent\textbf{Subcase 1: Two interior points.}
If $z_1,z_2\in B^c$ and $\overline{\xi}(z_1)=\overline{\xi}(z_2)$, then
$\overline{\xi}_{B^c}(z_1)=\overline{\xi}_{B^c}(z_2)$. By the univalence of $\overline{\xi}_{B^c}$ on
$B^c$, we get $z_1=z_2$.
\smallskip

\noindent\textbf{Subcase 2: One interior point, one branch cut point.}
Suppose $z\in B^c$ and $x\in B\setminus\{-1/\mu_n\}$.
The image $\overline{\xi}(B^c)=\overline{\xi}_{B^c}(B^c)$ is open.
The value $\overline{\xi}_B(x)$ is a one-sided (non-tangential) limit of $\overline{\xi}_{B^c}$
at $x\in\partial B^c$,
hence $\overline{\xi}_B(x)\in\partial\,\overline{\xi}_{B^c}(B^c)$ and is not an interior point.
Therefore $\overline{\xi}(z)\neq \overline{\xi}(x)$.
\smallskip

\noindent\textbf{Subcase 3: Two branch cut points.}
Let $x_1\neq x_2$ in $B\setminus\{-1/\mu_n\}$. Since $\overline{\xi}(x)=\overline{\xi}_B(x)$ on
$B\setminus\{-1/\mu_n\}$, it suffices to show $\overline{\xi}_B(x_1)\neq \overline{\xi}_B(x_2)$.

Suppose $x_1,x_2$ lie in one of the same half-intervals in \eqref{eff}.
Then, $\operatorname{Im} \overline{\xi}_B$ is constant while $\gamma_B'(x)$ has a fixed nonzero sign;
 hence $\overline{\xi}_B$ is strictly monotone along a horizontal line and
$\overline{\xi}_B(x_1)\neq \overline{\xi}_B(x_2)$.

On the other hand, if $x_1,x_2$ are in the two different halves of the same $I_n$,
then $\operatorname{Im} \overline{\xi}_B(x_1)=-\omega_n$ and $\operatorname{Im} \overline{\xi}_B(x_2)=+\omega_n$, so
$\overline{\xi}_B(x_1)\neq \overline{\xi}_B(x_2)$.

If $x_1,x_2$ lie in distinct intervals $I_j\neq I_n$, then
$\operatorname{Im} \overline{\xi}_B(x_\ell)\in\{\pm \omega_j,\pm \omega_n\}$ with $\omega_j\neq \omega_n$
since the sequence $\{\omega_n\}$ is strictly increasing in $n$. Again, the imaginary
parts differ, so the values are distinct.

Combining the three cases shows $\overline{\xi}$ is injective on
$\mathbb{C}_{\mu(\mathbb{T})} = \mathbb{C}\setminus\{-1/\mu_n\}$ when $M$ is finite.
\smallskip

\noindent\textbf{Case 2: $M$ countably infinite.}
If $M$ is countably infinite, by Lemma \ref{lem:B_structure},
 $$
 B\setminus \left\{-\frac{1} {\mu_n}\right\} = \bigcup_{n=1}^\infty I_n.
 $$
Since $I_0 = \emptyset$, we can define this simpler version of $\overline{\xi}_B:B\setminus\{-1/\mu_n\}\to\mathbb{C}$ by
\begin{equation}
 \overline{\xi}_B(x)=
 \begin{cases}
 \gamma_B(x)- i\,\omega_n, & x\in\big(-\tfrac{1}{\mu_{n+1}},\,\alpha_n\big],\\[2pt]
 \gamma_B(x)+ i\,\omega_n, & x\in\big(\alpha_n,\,-\tfrac{1}{\mu_n}\big),
 \end{cases}
\end{equation}
for all $n$ for which $\mu_{n+1}$ exists. Then
 \[
 \overline{\xi}(z)=
 \begin{cases}
 \overline{\xi}_{B^c}(z), & z\in B^c,\\[2pt]
 \overline{\xi}_B(x), & z=x\in B\setminus\{-1/\mu_n\}.
 \end{cases}
\]
The argument that $\overline{\xi}$ is injective in this situation is already contained in
the previous argument.
\end{proof}


The construction of the ergodic cylinder transformation in
Theorem \ref{thm:GlobalUnivalence}
gives us a way to extend the definition of $\gamma_{B^c}$ and $\omega_{B^c}$
to the branch cut via the functions $\gamma_B$ and $\omega_B$. We can tie them into
 global functions via
 $$
 \gamma(x,y) := \myRe(\overline{\xi}(x+iy),\quad\omega(x,y) := \myIm(\overline{\xi}(x+iy)),\quad\text{ for
 all }x+iy\in\mathbb{C}_{\mu(\mathbb{T})}.
 $$

Let $x+iy\in\mathbb{C}_{\mu(\mathbb{T})}$ be fixed. By Theorem~\ref{thm:GlobalUnivalence},
 $\overline{\xi}^{-1}: \overline{\xi}(\mathbb{C}_{\mu(\mathbb{T})}) \to \mathbb{C}_{\mu(\mathbb{T})}$ exists and is defined by
\[
\overline{\xi}^{-1}(\gamma+i \omega) = x+iy.
\]
where $\gamma(x,y)=\gamma$ and $\omega(x,y) = \omega$. The level curves of
$\gamma(x,y)$ and $\omega(x,y)$ therefore form a coordinate system on $\mathbb{C}_{\mu(\mathbb{T})}$
which is an orthogonal curvilinear coordinate system on $B^c$ by
Theorem \ref{thm:analytic}. See Figures~\ref{fig:erg_mapping3} and
\ref{fig:erg_mapping4}.

\begin{figure}[t]
  \centering
  \includegraphics[width=\textwidth]{fig_erg_mapping_3.pdf}
  \caption{The ergodic cylinder transformation on $\T=\Z$. Here, $\overline{\xi}^{-1}$ maps lines of constant average exponential growth rate (blue) and constant average exponential frequency (red) in the $\gamma$-$\omega$ plane to the level curves of $\gamma(x,y)$ (blue) and $\omega(x,y)$ (red) in the $x$-$y$ plane to form a new a curvilinear coordinate system on $\C$. We call this the ergodic complex plane for $\T$. Note that $\overline{\xi}^{-1}$ maps the imaginary axis in $\gamma$-$\omega$ coordinates to $\partial \mathcal{S}(\T)$ in the ergodic complex plane.}
  \label{fig:erg_mapping3}
\end{figure}

\begin{figure}[t]
  \centering
  \includegraphics[width=\textwidth]{fig_erg_mapping_4.pdf}
  \caption{The ergodic cylinder transformation for $\mathbb{T}_{1,5,0.45}$. The solid black curves are where $\gamma(x,y)=0$, i.e. $\partial \mathcal{S}(\mathbb{T})$.}
  \label{fig:erg_mapping4}
\end{figure}


\begin{definition} \rm
 Let $\mathbb{T}$ be a time scale. The regressive complex plane $\mathbb{C}_{\mu(\mathbb{T})}$ equipped with
 the $(\gamma,\omega)$ coordinate system induced by $\overline{\xi}^{-1}$ is called
the \emph{ergodic complex plane} for $\mathbb{T}$.
\end{definition}

\begin{theorem}\label{thm:regions}
Suppose \eqref{P} holds. Let $\mathbb{T}$ be a time scale which is unbounded above but with
bounded graininess. Let
 \begin{align*}
 \Gamma_0&:=\{(x,y):\gamma(x,y)=0\},\\
 \Gamma_{-}&:=\{(x,y):\gamma(x,y)<0\},\\
 \Gamma_{+}&:=\{(x,y):\gamma(x,y)>0\}.
 \end{align*}
\pagebreak

Then:
 \begin{enumerate}
 \item $\Gamma_0=\partial\mathcal{S}(\mathbb{T})$, $\Gamma_{-}$ is the interior of $\mathcal{S}(\mathbb{T})$, and $\Gamma_{+}$ is the exterior of $\mathcal{S}(\mathbb{T})$.
 \item $\overline{\xi}(\Gamma_0)=\{(0,\omega):-\Omega < \omega \leq \Omega\}$,
 $\overline{\xi}(\Gamma_{-})=\{(\gamma,\omega):\gamma<0\}$, and
 $\overline{\xi}(\Gamma_{+})=\{(\gamma,\omega):\gamma>0\}$.
 \item Let $\Gamma_{c}:=\{(x,y):-1/\mu_\textup{max}<x<0,y=0\}$. Then $\overline{\xi}(\Gamma_c)=\mathbb{R}^-$.
 \end{enumerate}
\end{theorem}

\begin{proof}
 (1) Assuming \eqref{P}, on $\Gamma_0$, \eqref{g21} vanishes. The second claim is
a restatement of \cite[Lemma~3.4b]{PoSiWi}. The third follows from the fact that
\textup{(C1)} is necessary for exponential stability, and \textup{(C1)} becomes
\eqref{pswtype} due to \eqref{P}.

(2) and (3) follow from the definitions of the sets involved and $\overline{\xi}$.
\end{proof}

Just as Hilger's pure imaginary number parameterizes the Hilger circle via
$$
\mathbb{I}_{\mu(t)} = \{ \stackrel{o}{\iota} \omega \mid -\pi/\mu(t)
< \omega \leq \pi/\mu(t) \},
$$
the ergodic cylinder transformation allows us to parameterize
$\partial \mathcal{S}=\Gamma_0$ as the preimage of the imaginary axis in
$\mathbb{C}^{\Omega}$.

\begin{definition} \rm
Suppose $-\Omega < \omega \leq \Omega$ such that $ \overline{\xi}^{-1}(i \omega)$ exists.
The \emph{ergodic pure imaginary number} $\overset{\odot}{\imath} \omega\in\partial\mathcal{S}=\Gamma_0$
is defined as
 \[ 
\overset{\odot}{\imath} \omega:= \overline{\xi}^{-1}(i \omega).
 \]
\end{definition}

On $\mathbb{R}$, $\overset{\odot}{\imath}\omega$ is $i\omega$ on the imaginary
axis,  and on $h\mathbb{Z}$, $\overset{\odot}{\imath}\omega$ is
 $\stackrel{o}{\iota} \omega$, the Hilger pure imaginary number, on
$\mathbb{I}_h$. This characterization allows us to parameterize $\Gamma_0$ as
\[
\Gamma_0 = \{ \overset{\odot}{\imath} \omega \mid -\Omega < \omega \leq \Omega \text{ and }
\overset{\odot}{\imath} \omega \text{ exists}\}.
\]
This parameterization allows us to prove properties of $\Gamma_0$, as the next two
propositions show.

\begin{proposition} %Lemma 2.6
Assume \eqref{P}. Let $\mathbb{T} = \mathbb{T}_{\mu_1,\mu_2,\ldots,\mu_n}$.
Then $\Gamma_0$ is bounded.
\end{proposition}

\begin{proof}
Fix $\omega \in (-\Omega,\Omega]$ such that $\overline{\xi}^{-1}(i \omega)$ exists and consider
the point $z=0+i\omega \in \mathbb{C}^{\Omega}$. Let $\mu_{\min}:=\min\mu_k$.
 Then $\overline{\xi}^{-1}(z)= \overset{\odot}{\imath} \omega\in\Gamma_0$, so
\begin{align*}
 0 & = \int_{0}^{L} \frac{\ln|1+\overset{\odot}{\imath} \omega \mu(t)|}{\mu(t)} \Delta t \\
 & \geq \int_{0}^{L} \frac{2 \myRe(\overset{\odot}{\imath} \omega) + |\overset{\odot}{\imath} \omega|^2 \mu(t)}{|1+\overset{\odot}{\imath} \omega \mu(t)|^2} \Delta t \\
 & \geq \int_{0}^{L} \frac{2 \myRe(\overset{\odot}{\imath} \omega) + |\overset{\odot}{\imath} \omega|^2 \mu_{\min}}{|1+\overset{\odot}{\imath} \omega \mu(t)|^2} \Delta t.
\end{align*}
Since the denominator is positive, this implies
$2 \myRe(\overset{\odot}{\imath} \omega) + |\overset{\odot}{\imath} \omega|^2 \mu_{\min} <0$, which means
$\overset{\odot}{\imath} \omega$ is in a disk of radius $1/\mu_{\min}$ centered at $-1/\mu_{\min}$
(which is a Hilger disk of curvature $\mu_{\min}$).
\end{proof}

\begin{proposition}
\label{prop:negCon}
$\Gamma_0$ is symmetric about the real axis. That is, for $\omega \in (-\Omega,\Omega]$
and $\overset{\odot}{\imath} \omega \not \in \mathbb{R}$, $\overline{\overset{\odot}{\imath} \omega} = \overset{\odot}{\imath} (-\omega)$.
\end{proposition}

\begin{proof}
Write $\overset{\odot}{\imath} \omega = x+iy$. Note $\overline{\xi}(x+iy) = \gamma(x,y) + i \omega(x,y) = i\omega$.
Then, calculations for $\overline{\overset{\odot}{\imath} \omega} = x-iy$ yield
\begin{align*}
 \gamma(x,-y)& =\frac{1}{L}\int_{0}^{L} \lim_{s \searrow \mu(t)}\frac{\ln |\overline{1+(x+iy)s}|}{s} \Delta s \\
 & = \frac{1}{L}\int_{0}^{L} \lim_{s \searrow \mu(t)}\frac{\ln |1+(x+iy)s|}{s} \Delta s \\
 & = \gamma(x,y)
  = 0,
\end{align*}
and
\begin{align*}
 \omega(x,-y)& =\frac{1}{L}\int_{0}^{L} \lim_{s \searrow \mu(t)}\frac{\text{Arg} (\overline{1+(x+iy)s})}{s} \Delta s \\
 & = - \frac{1}{L}\int_{0}^{L} \lim_{s \searrow \mu(t)}\frac{\text{Arg} (1+(x+iy)s)}{s} \Delta s \\
 & = -\omega(x,y)\\
 & = -\omega
\end{align*}
Therefore, $\overset{\odot}{\imath} (-\omega) = x-iy = \overline{\overset{\odot}{\imath} \omega}$.
\end{proof}

Finally, we now have a way of characterizing the time scale exponential function
in terms of its ergodic growth rate and ergodic frequency via the ergodic cylinder
transformation and its inverse.

\begin{remark}
\label{rmk:bigDeal}  \rm
Assuming \eqref{P}, an exponential function on $\mathbb{T}$ with ergodic growth rate
$\gamma$ and ergodic frequency $\omega$ and with constant subscript is given by
 \begin{equation}\label{exi}
 e_{\overline{\xi}^{-1}(\gamma+i \omega)}(t,t_0),
 \end{equation}
provided that $\overline{\xi}^{-1}(\gamma+i \omega)$ exists. In fact, among the family of functions of the form $e_\alpha(t,t_0)$ where $\alpha\in\mathbb{C}$ is constant, the
formulation in \eqref{exi} allows us to fully understand the qualitative behavior of $e_\alpha(t,t_0)$ in terms of the real and imaginary parts of $\alpha$.
\end{remark}

\begin{example}  \rm \quad

(1) The ergodic complex plane for $\mathbb{T}=\mathbb{T}_{1,2}$, the discrete periodic
time scale with graininesses $\{1,2\}$, is shown in Figure~\ref{fig:erg_plane2}. On this
time scale, $\Omega=2 \pi/3$.

(2) The ergodic complex plane for $\mathbb{T}=\mathbb{T}_{1,4,4}$, the discrete periodic
time scale with graininesses $\{1,4,4\}$, is shown in Figure~\ref{fig:erg_plane2}.
On this time scale, $\Omega=\pi/3$. In particular, the boundary of the region of
exponential stability on $\mathbb{T}_{1,4,4}$ contains a saddle point. The value of
$\omega$ at the saddle point is $2 \pi/9$. Therefore, $\overset{\odot}{\imath} 2\pi/9$ exists but
 $\overset{\odot}{\imath} (-2\pi/9)$ does not exist.

(3) The ergodic complex plane for $\mathbb{T}=\mathbb{P}_{1/4,1}$ is shown in
Figure~\ref{fig:erg_plane6}. Note that the boundary of the region of exponential
stability on $\mathbb{P}_{1/4,1}$ is disconnected. On the real axis, $\overset{\odot}{\imath} 4\pi/5$
is the first intersection of the boundary with $\mathbb{R}^-$, while $\overset{\odot}{\imath} (-4\pi/5)$ is
the second intersection of the boundary with $\mathbb{R}^-$, illustrating how we define
$\omega$ on $\mathbb{R}^-$. On this time scale, $\Omega=\infty$.

(4) Consider $\mathbb{T}=C_{1/3}$, the time scale consisting of repeated middle-third
Cantor sets. The ergodic complex plane for $\mathbb{T}=C_{1/3}$ is shown in
Figure~\ref{fig:erg_plane6}. On this time scale, $\Omega=\infty$.
\end{example}

\begin{example}
\label{ex:T12notBijective} \rm
The map $\overline{\xi}$ is in general not bijective. To see this, consider $\mathbb{T} = \mathbb{T}_{1,2}$.
The only saddle point of $\overline{\xi}_B$ is given by $\alpha_1 = -3/4$. Since we make the
choice to map saddle points of $\overline{\xi}_B$ to positive frequencies, we have
$\overline{\xi}(\alpha_1) = -\ln(2) + i \pi/3$. But, this means that the point
$z^{*} = -\ln 2-i\pi /3 \in \mathbb{C}^{\Omega}$ has no preimage under $\overline{\xi}$, that is,
$\overline{\xi}^{-1}(z^*)$ does not exist, and hence $\overline{\xi}$ is not a bijection between
$\mathbb{C}_{\mu(\mathbb{T})}$ and $\mathbb{C}^{\Omega}$.

In general, there are as many points in $\mathbb{C}_{\mu(\mathbb{T})}$ that do not have an image under
$\overline{\xi}$ as there are saddle points of $\overline{\xi}$, so for the time scales considered
in this paper, the number of such points is finite, and each of these points
is isolated.
\end{example}

We conclude this section by considering discrete time scales, which are prominent
 in time series data. We show that our results must adhere to major results in sampling
theory.

\begin{proposition}
Assume \eqref{P}. Let $\mathbb{T}=\mathbb{T}_{\mu_1,\mu_2,\ldots,\mu_N}$.
Then $L = \sum_{n=1}^{N} \mu_n$. Set
$a := \frac{L}{N}$. Then for $x+iy \in \mathbb{C}_{\mu(t)}$, $-\pi/a < \omega(x,y) \leq \pi/a$.
That is, $\Omega = \pi/a$.
\end{proposition}

\begin{proof}
Since $\mathbb{T}$ is discrete, for $x+iy \in B^c$,
$$
\omega_{B^c}(x,y)
= \frac{1}{L}\int_{0}^{L} \myIm_{\mu(t)}(x+iy) \Delta t
= \frac{1}{L}\sum_{k=1}^{N} \text{Arg}(1+(x+iy) \mu_k)
\leq \frac{N \pi}{L}= \frac{\pi}{a}.
$$
Moreover, on $B$, for $x \in I_0$ (that is, for $x$ to the left of all nonregressive
points), $\omega_B(x)=\pi/a$. For $x \in I_k, k \geq 1$, by construction
$\omega_B(x)<\pi/a$. Thus $\omega(x,y) \leq \pi/a$. The proof to show
$\omega(x,y)>-\pi/a$ is similar.
\end{proof}

\begin{remark} \rm
 If $\mathbb{T}=\mathbb{T}_{\mu_1,\mu_2,\ldots,\mu_N}$ and $t \in \mathbb{T}$ is measured
in seconds, then we have the following:

\textup{(1)} The average sampling rate is $1/a$ samples per second.

\textup{(2)} The Nyquist-Shannon Sampling Theorem \cite{Sh} says that we can perfectly
reconstruct a continuous time signal sampled on the time scale that contains no
frequencies higher than $\frac{1}{2a}$ hertz. That is, we can perfectly reconstruct
a continuous time signal sampled on the time scale that contains no frequencies higher
than $\pi/a = \Omega$, so $\Omega$ is the Nyquist frequency for sampling of a real
signal done on time scale points.

\textup{(3)} Because of (2), the definition of the ergodic frequency captures a key feature of
frequency from the perspective of signal processing.
\end{remark}

\begin{figure}[t]
  \centering
  \includegraphics[width=\textwidth]{fig_erg_plane_2.pdf}
\caption{The ergodic complex plane for $\T_{1,2}$ (left) and $\T_{1,4,4}$ (right). The solid black curve is where $\gamma(x,y)=0$, i.e. $\partial \mathcal{S}(\T)$.}
  \label{fig:erg_plane2}
\end{figure}


\begin{figure}[t]
  \centering
  \includegraphics[width=\textwidth]{fig_erg_plane_6.pdf}
  \caption{The ergodic complex plane for $\T=\mathbb{P}_{1/4,1}$ (left) and $\T=C_{1/3}$ (right).}
  \label{fig:erg_plane6}
\end{figure}


\section{The box plus operation}

In this section, we propose a generalization of the time scale circle plus binary
operator for the ergodic complex plane.

The cylinder strip $\mathbb{C}^{\Omega(h)}$ is a manifestation of $\mathbb{C}^{\Omega}$ when
$\mathbb{T} = h \mathbb{Z}$, where $h$ is constant, and therefore its coordinates can be thought of as a \emph{local}
growth rate on the real axis and \emph{local} frequency on the imaginary axis.
Therefore, by \eqref{eq:oplusViaXi}, one way to understand the action of $\oplus_h$
is that it maps $z,w\in\mathbb{C}_h$ to $\mathbb{C}^\Omega$ for $\mathbb{T}
=h\mathbb{Z}$, adds the local growth
rate of $z$ to the local growth rate of $w$ and adds the local frequency of $z$ to the
local frequency of $w \pmod{2\pi i/h}$, and finally pulls back to the unique point
in $\mathbb{C}_h$ that would map to that new growth rate and frequency pair under $\xi_h$.

We can generalize $\oplus_h$ by employing $\overline{\xi}$ rather than $\xi_{h}$ in
\eqref{eq:oplusViaXi}. Since $\overline{\xi}$ is not a bijection (see Example
\ref{ex:T12notBijective}), addition (mod $2 i \Omega$) in $\mathbb{C}_\Omega$ is not a group
since the sum of two points may not correspond to the image of a point under $\overline{\xi}$.
Thus $(\mathbb{C}^{\Omega},+\pmod{2 i \Omega} )$ is a groupoid, and $+\pmod{2 i \Omega}$
is a partial function.

\begin{example}
\label{ex:T12prob} \rm
To demonstrate the lack of closure under addition (mod $2 i \Omega$) in $\mathbb{C}^{\Omega}$,
consider again the time scale $\mathbb{T}_{1,2}$ from Example \ref{ex:T12notBijective}.

Fixing $z \in \mathbb{C}^{\Omega}$, there is just one point $v$ such that $z + v$ does not exist.
This point $v$ satisfies
\[
v = -\ln2-i\pi/3-z.
\]

There are a few properties we want to point out in this example that are illustrative
for later proofs. Notice that
$\overline{\xi}^{-1}(\overline{\xi}(-3/4)) = \overline{\xi}^{-1}(-\ln2+\pi i/3) = -3/4$.
Conversely, $\overline{\xi}(\overline{\xi}^{-1}(-\ln2-\pi i/3))$ does not exist. Finally,
$\overline{\xi}(\overline{\xi}^{-1}(\overline{\xi}(-3/4))) = \overline{\xi}(-3/4) = -\ln2+\pi i/3$.
This example gives us intuition for which compositions of $\overline{\xi}$ and $\overline{\xi}^{-1}$
exist in general.
\end{example}

Crucially, $\overline{\xi}^{-1}(\overline{\xi}(z))$ always exists, but $\overline{\xi}(\overline{\xi}^{-1}(z))$ may not
exist. Moreover,
$$
\overline{\xi}(\overline{\xi}^{-1}(\overline{\xi}(z))) = \overline{\xi}(z).
$$

Since $\overline{\xi}$ is defined as an integral over the time scale, $\overline{\xi}$ can act on
functions as well so long as the resulting integral converges. In what follows,
for fixed $z,w$, we will evaluate $\overline{\xi}$ at $z \oplus_{\mu(t)} w$, which is a
function of $t$.

\begin{proposition}
\label{thm:oplusXiBar}
Let $z,w \in B^c$.
Then
 \begin{gather*}
 \overline{\xi}_{B^c}(z \oplus_{\mu(t)} w) = \overline{\xi}_{B^c}(z) + \overline{\xi}_{B^c}(w) \pmod{ 2i \Omega},\\%\label{eq:oplusXiBar}\\
 \overline{\xi}_{B^c}(z \ominus_{\mu(t)} w) = \overline{\xi}_{B^c}(z) - \overline{\xi}_{B^c}(w) \pmod{2 i \Omega}. %\label{eq:ominusXiBar}
 \end{gather*}
\end{proposition}

\begin{proof}
Let $z,w \in \mathbb{C}_{\mu(\mathbb{T})}$. Then $z,w$ are regressive for each $t \in \mathbb{T}$, so
$z \oplus_{\mu(t)} w$ is defined for each $t \in \mathbb{T}$. Hence,
 \begin{align*}
 \overline{\xi}_{B^c}(z \oplus_{\mu(t)} w)
 & := \frac{1}{L} \int_{0}^{L} \xi_{\mu(t)}(z \oplus_{\mu(t)} w)\,\Delta t \pmod{2 i \Omega}\\
 & = \frac{1}{L}\int_{0}^{L} \xi_{\mu(t)}(z) + \xi_{\mu(t)}(w)\,\Delta t \pmod{2 i \Omega} \\
 & = \frac{1}{L}\int_{0}^{L} \xi_{\mu(t)}(z) \,\Delta t + \frac{1}{L}\int_{0}^{L} \xi_{\mu(t)}(w)\,\Delta t \pmod{2 i \Omega}\\
 & = \overline{\xi}_{B^c}(z) + \overline{\xi}_{B^c}(w) \pmod{2 i \Omega}.
 \end{align*}
The proof for $\ominus_{\mu(t)}$ follows similarly.
\end{proof}

Although according to Proposition~\ref{thm:oplusXiBar}, the function
$\overline{\xi}_{B^c}$ appears to be operation preserving, it is important to realize that
$z \oplus_{\mu(t)} w$ is not a point in $\mathbb{C}_{\mu(\mathbb{T})}$, but instead is a function of $t$.
Nevertheless, if $\overline{\xi}^{-1}(\overline{\xi}(z) + \overline{\xi}(w) \pmod{2 i \Omega})$ exists, there
is a unique point in $\mathbb{C}_{\mu(\mathbb{T})}$ that maps to
$\overline{\xi}(z) + \overline{\xi}(w) \pmod {2 i \Omega} $. This fact, along with
Proposition~\ref{thm:oplusXiBar}, motivates us to call this point the \emph{box plus
sum of $z$ and $w$}, denoted $z \boxplus_{\mathbb{T}} w$. We can similarly define the
\emph{box minus difference of $z$ and $w$}, denoted $z \boxminus_{\mathbb{T}} w$.

\begin{definition} \rm
Let $z,w \in \mathbb{C}_{\mu(\mathbb{T})}$ and suppose
$\overline{\xi}^{-1}(\overline{\xi}(z) + \overline{\xi}(w)\pmod {2 i \Omega} )$ exists.
We define the partial function
$\boxplus_{\mathbb{T}}:\mathbb{C}_{\mu(\mathbb{T})} \times \mathbb{C}_{\mu(\mathbb{T})} \to \mathbb{C}_{\mu(\mathbb{T})}$ by
\[ %\label{eq:boxPlusDef}
 z \boxplus_{\mathbb{T}} w := \overline{\xi}^{-1}(\overline{\xi}(z) + \overline{\xi}(w) \pmod {2 i \Omega}  ).
\]
Similarly, if $z,w \in \mathbb{C}_{\mu(\mathbb{T})}$ and $\overline{\xi}^{-1}(\overline{\xi}(z)
- \overline{\xi}(w) \pmod{2 i \Omega} )$ exists, then define the partial function
$\boxminus_{\mathbb{T}}:\mathbb{C}_{\mu(\mathbb{T})} \times \mathbb{C}_{\mu(\mathbb{T})} \to \mathbb{C}_{\mu(\mathbb{T})}$ by
\[%\label{eq:boxMinusDef}
 z \boxminus_{\mathbb{T}} w := \overline{\xi}^{-1}(\overline{\xi}(z) - \overline{\xi}(w) \pmod{2 i \Omega}  ).
\]
Finally, if $z \in \mathbb{C}_{\mu(\mathbb{T})}$ and $\overline{\xi}^{-1}(- \overline{\xi}(w) )$ exists, define
 $\boxminus_{\mathbb{T}} z := 0 \boxminus_{\mathbb{T}} z$. Note the mod condition is not necessary
in the last definition since $\overline{\xi}(w) \in \mathbb{C}^{\Omega}$ implies
$-\overline{\xi}(w) \in \mathbb{C}^{\Omega}$.
\end{definition}

Now, we have the following useful identities.

\begin{proposition}
\label{prop:boxPlusList}
Let $z,w,u \in \mathbb{C}_{\mu(\mathbb{T})}$. Then
\begin{enumerate}
 \item $0 \boxplus_{\mathbb{T}} z = z$
 \item If $z \boxplus_{\mathbb{T}} w$, $w \boxplus_{\mathbb{T}} u$, $(z \boxplus_{\mathbb{T}} w) \boxplus_{\mathbb{T}} u$, and $z \boxplus_{\mathbb{T}} (w \boxplus_{\mathbb{T}} u)$ exist, then $$(z \boxplus_{\mathbb{T}} w) \boxplus_{\mathbb{T}} u = z \boxplus_{\mathbb{T}} (w \boxplus_{\mathbb{T}} u).$$
 \item If $z \boxplus_{\mathbb{T}} w$ exists, then $z \boxplus_{\mathbb{T}} w = w \boxplus_{\mathbb{T}} z$.
 \item If $(\boxminus_{\mathbb{T}} z)$ exists, then $z \boxminus_{\mathbb{T}} z = z \boxplus_{\mathbb{T}} (\boxminus_{\mathbb{T}} z) = 0$.
 \item If $(\boxminus_{\mathbb{T}} z)$ exists, then $\boxminus_{\mathbb{T}}(\boxminus_{\mathbb{T}} z) = z$.
 \item If $z \boxminus_{\mathbb{T}} w$ and $w \boxminus_{\mathbb{T}} z$ exist, then $\boxminus_{\mathbb{T}}(z \boxminus_{\mathbb{T}} w) = w \boxminus_{\mathbb{T}} z$.
 \item If $(\boxminus_{\mathbb{T}} z)$ and $(\boxminus_{\mathbb{T}} w)$ exist, then $\boxminus_{\mathbb{T}}(z \boxplus_{\mathbb{T}} w) = (\boxminus_{\mathbb{T}} z) \boxplus_{\mathbb{T}} (\boxminus_{\mathbb{T}} w)$.
 \item If $z \in \Gamma_0 / \mathbb{R}$, then $\boxminus_{\mathbb{T}} z$ exists and
 $\overline{z} = \boxminus_{\mathbb{T}} z$. Thus, if $z = \overset{\odot}{\imath} \omega$, then $\overset{\odot}{\imath}(-\omega) = \boxminus_{\mathbb{T}} \overset{\odot}{\imath} \omega$.
 \end{enumerate}
 \end{proposition}

We will prove (1), (2), (5), and (8). The proofs of (3), (4), (6) and (7) are
similar to the proof of (2), in the sense the proof involves compositions and
cancelations of $\overline{\xi}$ and $\overline{\xi}^{-1}$.

\begin{proof}
(1) $0 \boxplus_{\mathbb{T}} z = \overline{\xi}^{-1}(\overline{\xi}(0) + \overline{\xi}(z) \pmod {2 i \Omega} ) 
= \overline{\xi}^{-1}( 0 + \overline{\xi}(z)) =\overline{\xi}^{-1}( \overline{\xi}(z)) = z$.

(2) Since $z \boxplus_{\mathbb{T}} w$ exists,
$\overline{\xi}(\overline{\xi}^{-1}(\overline{\xi}(z) + \overline{\xi}(w) \pmod{2 i \Omega}))
 = \overline{\xi}(z) + \overline{\xi}(w) \pmod{2 i \Omega}$. A similar expression holds
for $w \boxplus_{\mathbb{T}} u$. Then
\begin{align*}
 (z \boxplus_{\mathbb{T}} w) \boxplus_{\mathbb{T}} u &= \overline{\xi}^{-1} ( \overline{\xi}(\overline{\xi}^{-1}(\overline{\xi}(z) + \overline{\xi}(w) \pmod {2 i \Omega})) + \overline{\xi}(u) \pmod{2 i \Omega}) \\
 & = \overline{\xi}^{-1} ( (\overline{\xi}(z) + \overline{\xi}(w) \pmod{2 i \Omega} ) + \overline{\xi}(u) \pmod{2 i \Omega} ) \\
 & = \overline{\xi}^{-1} ( \overline{\xi}(z) + (\overline{\xi}(w)+ \overline{\xi}(u) \pmod {2 i \Omega}) \pmod {2 i \Omega}) \\
 & = \overline{\xi}^{-1} ( \overline{\xi}(z) + \overline{\xi}(\overline{\xi}^{-1}(\overline{\xi}(w) + \overline{\xi}(u) \pmod {2 i \Omega}))\pmod {2 i \Omega}) \\
 & = z \boxplus_{\mathbb{T}} (w \boxplus_{\mathbb{T}} u).
\end{align*}

(5) Consider
 \begin{align*}
 \boxminus_{\mathbb{T}}(\boxminus_{\mathbb{T}} z) = \boxminus_{\mathbb{T}}(\overline{\xi}^{-1}(-\overline{\xi}(z))
 = \overline{\xi}^{-1}(-\overline{\xi}(\overline{\xi}^{-1}(-\overline{\xi}(z))))
 = \overline{\xi}^{-1}(\overline{\xi}(z))= z.
 \end{align*}

(8) Since $z \in \Gamma_0/\mathbb{R}$, there is a unique $\omega \in (-\Omega,\Omega]$
such that $z = \overline{\xi}^{-1}(i \omega) = \overset{\odot}{\imath} \omega$. Then
 \begin{align*}
 \boxminus_{\mathbb{T}} z = \boxminus_{\mathbb{T}} \overset{\odot}{\imath} \omega
 = \overline{\xi}^{-1}(-\overline{\xi}(\overline{\xi}^{-1}(i \omega)))
 = \overline{\xi}^{-1}(-i \omega)
 = \overset{\odot}{\imath} (-\omega)
 = \overline{z},
 \end{align*}
where we use Proposition \ref{prop:negCon} for the last equality.
\end{proof}

Thus, $\boxplus_\mathbb{T}$ provides a global analogue of the local action of $\oplus_{\mu(t)}$.
The $\boxplus_{\mathbb{T}}$ operator maps $z,w\in\mathbb{C}_\mu$ to $\mathbb{C}_\Omega$, adds the ergodic
growth rate of $z$ to the ergodic growth rate of $w$ and adds the ergodic frequency
of $z$ to the ergodic frequency of $w$ (mod $2i \Omega$), and finally pulls back to
the unique point (if it exists) in $\mathbb{C}_{\mu(\mathbb{T})}$ that would map to that new growth
rate and frequency pair under $\overline{\xi}$.

The operations $\boxminus_\mathbb{T}$ and $\ominus_\mu$ also interact with time scale
exponential functions in an interesting way. To show that, we need the following
result.

\begin{proposition}
Let $\mathbb{T}$ be a time scale and let $-\Omega < \omega \leq \Omega$ such that
$\overset{\odot}{\imath} \omega$ exists. When $\overset{\odot}{\imath}(-\omega) = \boxminus_{\mathbb{T}} (\overset{\odot}{\imath} \omega)$
exists, we have
$$
\overline{e_{\overset{\odot}{\imath} \omega}(t,t_0)} = e_{\boxminus \overset{\odot}{\imath} \omega}(t,t_0).
$$
\end{proposition}

\begin{proof}
(1) This follows since $\overline{e_{z}(t,t_0)} = e_{\overline{z}}(t,t_0)$ for
regressive $z$ and since $\overset{\odot}{\imath} (-\omega) = \overline{\overset{\odot}{\imath} \omega} = \boxminus_{\mathbb{T}}
 \overset{\odot}{\imath} \omega$ by Propositions \ref{prop:negCon} and \ref{prop:boxPlusList}.
\end{proof}

We now can show the relationship between $\ominus_{\mu(t)}$ and $\boxminus_\mathbb{T}$ in the
time scale exponential. Since $\ominus_{\mu(t)}$ is a \emph{local} operation while
$\boxplus_{\mathbb{T}}$ is a \emph{global} operation, the following proposition provides a
bridge between the local and the global behavior of the time scale exponential
function.

\begin{proposition}
\label{prop:oplusBoxPlus}
Let $\mathbb{T}$ be a time scale and let $-\Omega < \omega \leq \Omega$ such that $\overset{\odot}{\imath} \omega$ exists. When $\boxminus_{\mathbb{T}} (\overset{\odot}{\imath} \omega)$ exists,
$$
e_{\ominus_{\mu(t)} \overset{\odot}{\imath} \omega}(t,t_0) = e_{\boxminus_{\mathbb{T}} \overset{\odot}{\imath} \omega}(t,t_0) Q_{\omega}(t,t_0),
$$
where $Q_\omega(t,t_0):=1/|e_{\overset{\odot}{\imath} \omega}(t,t_0)|^2$.
\end{proposition}

\begin{proof}
By direct calculation,
 $$
 e_{\ominus_{\mu(t)} \overset{\odot}{\imath} \omega}(t,t_0) = \frac{1}{e_{\overset{\odot}{\imath} \omega}(t,t_0)}
 = \frac{\overline{e_{\overset{\odot}{\imath} \omega}(t,t_0)}}{|e_{\overset{\odot}{\imath} \omega}(t,t_0)|^2}
 = e_{\boxminus_{\mathbb{T}} \overset{\odot}{\imath} \omega}(t,t_0) Q_{\omega}(t,t_0).
 $$
\end{proof}

Notice in Proposition \ref{prop:oplusBoxPlus} that $\boxminus_{\mathbb{T}} \overset{\odot}{\imath} \omega$ is a
number, while $\ominus_{\mu(t)} \overset{\odot}{\imath} \omega$ is a function of $t$.
The effect that the time variation in $\ominus_{\mu(t)} \overset{\odot}{\imath} \omega$ has on the
time scale exponential is embedded in $Q_{\omega}(t,t_0)$.

As explained in Remark \ref{rmk:bigDeal}, each point $z \in \mathbb{C}_{\mu(\mathbb{T})}$
has an ergodic growth rate and ergodic frequency. Using $\boxplus_{\mathbb{T}}$, we can
explicitly write $z$ in an analogue of the rectangular form, much like the Hilger
decomposition.

Generalizing the Hilger real and imaginary parts in the spirit of Lemma~\ref{lem:HilXi},
we have the following definition.

\begin{definition} \rm
Let $z \in \mathbb{C}_{\mu(\mathbb{T})}$. We define the \emph{ergodic real part} of $z$, denoted
$\myRe_{\mathbb{T}}(z)$, by
 \[ %\label{eq:ergRealPart}
 \myRe_{\mathbb{T}}(z):=\overline{\xi}^{-1}(\myRe(\overline{\xi}(z))),
 \]
and the \emph{ergodic imaginary part} of $z$, denoted $\myIm_{\mathbb{T}}(z)$, by
\[
 %\label{eq:ergImagPart}
\myIm_{\mathbb{T}}(z):=\myIm(\overline{\xi}(z)).
\]\end{definition}

Note that $\myRe_{\mathbb{T}}(z)$ is the right-most point on the real axis along the contour
$\gamma = \overline{\xi}(z)$ by virtue of having zero ergodic imaginary part. In particular,
this point is to the right of $-1/\mu_{\max}$, which means $\overline{\xi}^{-1}(\myRe(\overline{\xi}(z)))$
always exists. Also, the ergodic growth rate $\gamma(x,y)$ and ergodic real part
$\myRe_{\mathbb{T}}(x+iy)$ are related by
$$
 \overline{\xi}(\myRe_{\mathbb{T}}(x+iy)) = \myRe(\overline{\xi}(x+iy)) = \gamma(x,y).
$$

We are now able to give a rectangular ergodic decomposition for $z \in \mathbb{C}_{\mu(\mathbb{T})}$
in terms of $\boxplus_\mathbb{T}$. See Figure~\ref{fig:ergodic_decomp}.

\begin{figure}[t]
 \centering
\includegraphics[scale=1.5]{fig_ergodic_decomp.pdf}
 \caption{The ergodic decomposition $z=\myRe_\mathbb{T} z\boxplus_\mathbb{T} \myIm_\mathbb{T} z$. $\myRe_\mathbb{T} z$ is found by projecting along the appropriate level curve of $\gamma(x,y)$ until it reaches the $x$-axis while $\myIm_\mathbb{T} z$ is found by projecting along the appropriate level curve of $\omega(x,y)$ until it reaches $\Gamma_0$. If $z\in\Gamma_{-}$, then $\myRe_\mathbb{T} z<0$, while if $z\in\Gamma_{+}$, then $\myRe_\mathbb{T} z>0$. If $\myIm z>0$, then $\myIm_\mathbb{T} z>0$, while if $\myIm z<0$, then $\myIm_\mathbb{T} z<0$. Care must be taken at self-crossings such as $c$, where our definition chooses $\myIm_\mathbb{T} c>0$.}
 \label{fig:ergodic_decomp}
\end{figure}

\begin{proposition}
Let $z \in \mathbb{C}_{\mu(\mathbb{T})}$. Then $z = \myRe_{\mathbb{T}}(z) \boxplus_{\mathbb{T}} \overset{\odot}{\imath} \myIm_{\mathbb{T}}(z)$.
\end{proposition}

\begin{proof}
\begin{align*}
\myRe_{\mathbb{T}}(z) \boxplus_{\mathbb{T}} \overset{\odot}{\imath} \myIm_{\mathbb{T}}(z)
& = \overline{\xi}^{-1}(\overline{\xi}(\myRe_{\mathbb{T}}(z)) + \overline{\xi}(\overset{\odot}{\imath} \myIm_{\mathbb{T}}(z)) \pmod {2 i \Omega }) \\
& = \overline{\xi}^{-1}(\myRe(\overline{\xi}(z)) + i \myIm_{\mathbb{T}}(z)\pmod {2 i \Omega}) \\
& = \overline{\xi}^{-1}(\myRe(\overline{\xi}(z)) + i \myIm(\overline{\xi}(z)) \pmod {2 i \Omega}) \\
& = \overline{\xi}^{-1}(\overline{\xi}(z))
 = z.
\end{align*}
\end{proof}

\begin{example}
 Again, consider $\mathbb{T}_{1,2}$. The rectangular ergodic decomposition exists even for
$\alpha_1=-3/4$. Indeed, by the calculations in Example~\ref{ex:T12prob},
 $ \myRe_{\mathbb{T}_{1,2}}(z) = \overline{\xi}^{-1}(-\ln 2) \approx -0.396447 $
and
 $ \myIm_{\mathbb{T}_{1,2}}(z) = \frac{\pi}{3}$.
Thus,
 \[
 \alpha_1 = \overline{\xi}^{-1}(-\ln(2)) \boxplus_{\mathbb{T}_{1,2}} \left( \overset{\odot}{\imath} \frac{\pi}{3} \right)
 \approx (-0.396447) \boxplus_{\mathbb{T}_{1,2}} (-0.75 + 0.661438 i).
 \]
\end{example}

\begin{remark}
\label{rmk:group} \rm
We can recover a group structure for the box plus operator by extending $\overline{\xi}$
so its codomain is all of $\mathbb{C}_{\mu(\mathbb{T})}$ by changing the domain so that a point is
also associated with its frequency. Consider the set
$\mathbb{C}_{\mu(\mathbb{T})}^g \subset \mathbb{C}_{\mu(\mathbb{T})} \times (-\Omega,\Omega]$ with the condition that
$(z,\omega)\in\mathbb{C}_{\mu(\mathbb{T})}^{g}$ if and only if $\myIm(\overline{\xi}(z))=\omega$ when $z$
is not a saddle point of $\overline{\xi}$ and $\myIm(\overline{\xi}(z)) = \pm \omega$ when $z$ is a
saddle point of $\overline{\xi}$. Then define $\overline{\xi}^g:\mathbb{C}_{\mu(\mathbb{T})}^g \to \mathbb{C}^{\Omega}$ by
\[
 \overline{\xi}^g(z,\omega) = \begin{cases}
\gamma + i \omega, \quad & \overline{\xi}(z)=\gamma+i\omega, \\
\gamma + i \omega, \quad & \overline{\xi}(z)=\gamma-i\omega.
 \end{cases}
\]
Then, for a saddle point $\alpha_k$ of $\overline{\xi}_B$ satisfying
$\overline{\xi}_B(\alpha_k)=\gamma+i\omega$,
\[
(\overline{\xi}^g)^{-1}(\gamma-i \omega) = (\alpha_k,-\omega),
\]
and for all other points $(\overline{\xi}^{g})^{-1}$ agrees with $\overline{\xi}^{-1}$.

With this, define
$\boxplus_{\mathbb{T}}^{g}:\mathbb{C}_{\mu(\mathbb{T})}^{g}\times\mathbb{C}_{\mu(\mathbb{T})}^{g}\to\mathbb{C}^{\Omega}$ by
\[%\label{eq:boxPlusDef}
 (z_1,\omega_1) \boxplus_{\mathbb{T}}^g (z_2,\omega_2)
 := (\overline{\xi}^{g})^{-1}(\overline{\xi}^g(z_1,\omega_1) + \overline{\xi}^g(z_2,\omega_2)),
\]
where the addition is $\text{mod } {2 i \Omega}$ in the first argument of $\overline{\xi}^g$ and is
$\text{mod } 2 \Omega$ in the second argument of $\overline{\xi}^g$.
Then $(\mathbb{C}_{\mu(\mathbb{T})}^g,\boxplus_{\mathbb{T}}^g)$ is a group with identity $(0,0)$ and with
inverse of $(z,\omega)$ being $((\overline{\xi}^g)^{-1}(-\overline{\xi}^g(z)),-\omega)$.
\end{remark}

\begin{remark} \rm
For an easier group structure involving box plus, we can generalize the observation
that the Hilger circle with the circle plus, $({\mathcal H}_{\mu(t)},\oplus_{\mu(t)})$,
is a subgroup of $(\mathbb{C}_{\mu(t)},\oplus_{\mu(t)})$ and is isomorphic to
$((-\pi/\mu(t),\pi/\mu(t)],+ \pmod {2 \pi/\mu(t)})$. Notice that as long as $\Gamma_0$
does not have any self-intersections, $\overline{\xi}:\Gamma_0 \to \overline{\xi}(\Gamma_0)$ is a
bijection. In this case, $(\Gamma_0,\boxplus_{\mathbb{T}})$ is a group which is isomorphic
to $((-\Omega,\Omega],+ \pmod{2 \Omega})$.
\end{remark}

\begin{remark} \rm
We can extend the range of $\omega$ from the Nyquist range $(-\Omega,\Omega]$ to $\mathbb{R}$.
Fix the positive principal branch of $\overline{\xi}^{-1}$ on the Nyquist segment
$i(-\Omega,\Omega]\setminus\{i\omega_{-}\}$, where $i\omega_{-}$ is the
parameter corresponding to a self-intersection point of $\Gamma_0$.
Analytic continuation along the imaginary axis lifts this
branch to the Riemann surface $\mathcal R_{\overline{\xi}^{-1}}$, giving a
single-valued map
\[
 \overline{\xi}^{-1}_{\textup{lift}}\colon i\bigl(\mathbb{R}\setminus\Lambda\bigr)
 \to\mathcal{C},
 \quad
 \Lambda:=\{\omega_{-}+2k\Omega: k\in\mathbb{Z}\},
\]
which satisfies
$\overline{\xi}^{-1}_{\textup{lift}}\bigl(i(\omega+2\Omega)\bigr)=
 \overline{\xi}^{-1}_{\textup{lift}}(i\omega)$ for every
$\omega\in\mathbb{R}\setminus\Lambda$.
Projecting to the plane extends $\overline{\xi}^{-1}$ to the punctured imaginary
axis while identifying all inputs that differ by $2\Omega$.
Let $Z=\operatorname{im}\overline{\xi}^{-1}_{\textup{lift}}$.
Because the branch on
$i(-\Omega,\Omega]\setminus\{i\omega_{-}\}$ is injective, every
$z\in Z$ has a unique representative
$\omega\in(-\Omega,\Omega]\setminus\{\omega_{-}\}$.
We define
\[
 z_1\boxplus z_2:=
 \overline{\xi}^{-1}_{\textup{lift}}\!\bigl(i(\omega_1+\omega_2)\bigr),
 \quad
 \text{whenever }\omega_1+\omega_2\notin\Lambda.
\]
Hence $\overline{\xi}^{-1}_{\textup{lift}}$ is an isomorphism onto $(Z,\boxplus)$.
On the punctured circle this provides a global operation; on all of
$\mathbb{R}\setminus 2\Omega\mathbb{Z}$, it is only partial because inputs whose sum is in $\Lambda$
are excluded. However, a construction similar to that in Remark~\ref{rmk:group}
will create a global group isomorphism.
\end{remark}

\section{Conclusion}
In this paper we introduced the ergodic complex plane as a global analogue of Hilger's
complex plane on time scales. By averaging the cylinder transformation, we defined
the ergodic growth rate and ergodic frequency, and showed that together they generate
the ergodic cylinder transformation, an analytic, globally univalent map that yields
an orthogonal curvilinear coordinate system on the regressive complex plane.
This framework provides an ergodic decomposition analogous to Hilger's as well as the
global box plus operation as an extension of the local circle plus operation.
These results unify local and global perspectives on growth and frequency, and open
the way for further developments in harmonic and spectral analysis on general time
scales.

\begin{thebibliography}{10}

\bibitem{allouche1999ubiquitous}
Jean-Paul Allouche, Jeffrey Shallit;
 \emph{The ubiquitous
  {P}rouhet-{T}hue-{M}orse sequence}, Sequences and their Applications:
  Proceedings of SETA’98, Springer, 1999, pp.~1--16.

\bibitem{AvAk}
FG~Avkhadiev, LA~Aksent'ev;
 \emph{The main results on sufficient conditions
  for an analytic function to be schlicht}, Russian Mathematical Surveys,
  \textbf{30} (1975), no.~4, 1.

\bibitem{BoPeD}
Martin Bohner, Allan Peterson;
\emph{Dynamic {E}quations on {T}ime {S}cales:
  {A}n {I}ntroduction with {A}pplications}, Birkh\"auser, 2001.

\bibitem{DaGrMaJa}
John~M Davis, Ian~A Gravagne, Robert~J Marks, Billy~J Jackson;
  \emph{Regions of exponential stability for {LTI} systems on nonuniform
  discrete domains}, IEEE Proceedings of the 43rd Southeastern Symposium on
  System Theory, Auburn, AL (2011), 37--42.

\bibitem{HiA}
Stefan Hilger;
\emph{Analysis on measure chains---a unified approach to
  continuous and discrete calculus}, Results in mathematics \textbf{18} (1990),
  no.~1, 18--56.

\bibitem{JaDaE}
Billy~J Jackson, John~M Davis, \emph{An ergodic approach to
  {P}\"otzsche-{S}iegmund-{W}irth's region of exponential stability}, Nonlinear
  Analysis: Hybrid Systems \textbf{36} (2020), 100848.

\bibitem{JaDalap}
Billy~J {J}ackson and John~M {D}avis;
\emph{An ergodic approach to {L}aplace
  transforms on time scales}, Journal of Mathematical Analysis and Applications
  \textbf{502} (2021), no.~1, 125231.

\bibitem{Ka}
Ba{\c{s}}ak Karpuz;
 \emph{Analyticity of the complex time scale exponential},
  Complex Analysis and Operator Theory, \textbf{11} (2017), 21--34.

\bibitem{PoSiWi}
Christian Potzsche, Stefan Siegmund, Fabian Wirth;
\emph{A spectral
  characterization of exponential stability for linear time-invariant systems
  on time scales}, Discrete and Continuous Dynamical Systems, \textbf{9} (2003),
  no.~5, 1223--1242.

\bibitem{PoDaGr}
Dylan Poulsen, John Davis, Ian Gravagne;
\emph{Mean square stability for
  systems on stochastically generated discrete time scales}, Nonlinear
  Analysis: Hybrid Systems, \textbf{31} (2019), 41--55.

\bibitem{Sh}
Claude~E Shannon;
 \emph{Communication in the presence of noise}, Proceedings of
  the IRE, \textbf{37} (2006), no.~1, 10--21.

\end{thebibliography}

\end{document}
